%%%%%%%%%%%%%%%%%%%%%%%%%%%%%%%%%%%%%%%%%%%%%%%%%%%%%%%%%%%%%
%% Begin exercise %%
%%%%%%%%%%%%%%%%%%%%%%%%%%%%%%%%%%%%%%%%%%%%%%%%%%%%%%%%%%%%%
\ex{Distributed and Concentrated Windings}


\normalsize{\textbf{Acknowledgement}: The following exercise is adapted from ``Geregelte Drehstromantriebe / Controlled AC Drives'' by J. Böcker, Paderborn University, 2021 and 
``Elektrische Maschinen und Antriebe Übungsbuch: Aufgaben mit Lösungsweg'' by A. Binder, Springer, 2017
}\\


%%%%%%%%%%%%%%%%%%%%%%%%%%%%%%%%%%%%%%%%%%%%%%%%%%%%%%%%%%%%%
%% Task 1 %%
%%%%%%%%%%%%%%%%%%%%%%%%%%%%%%%%%%%%%%%%%%%%%%%%%%%%%%%%%%%%%

\task{Distributed windings}
Consider a 4-pole three-phase motor with 15 stator slots. A simplified sketch of this motor with the winding scheme of phase a is shown in \autoref{fig:MMF_distributed}.

\begin{figure}[htb]
    \centering
    \includegraphics{ex04/MMF_distributed.pdf}
    \caption{Simplified sketch of a distributed winding scheme. Only phase a is shown.}
    \label{fig:MMF_distributed}
\end{figure}



%%%%%%%%%%%%%%%%%%%%%%%%%%%%%%%%%%%%%%%%%%%%%%%%%%%%%%%%%%%%%
\subtask{Determine the pole pitch $\rho_{\mathrm{p}}$ and the number of notches. Is it an integral-slot or a fractional-slot winding?}

\begin{solutionblock}
    The number of slots are shown in the figure, which leads to $Q$ = 15. Furthermore, the number of phases is given in the task, that is $m$ = 3.

    Hence, the pole pitch is calculated as:
    \begin{equation}
        \rho_{\mathrm{p}} = \frac{2\pi}{2p}
        = \frac{\pi}{2}.
    \end{equation}
    
    The number of notches is defined with:
    \begin{equation}
        q = \frac{Q}{2pm} = \frac{15}{2\cdot 2 \cdot 3}
        = \frac{15}{12} = \frac{5}{4}.
    \end{equation}
    Since the number of notches per pole and phase is not an integer, it is a fractional slot winding.
    The coil width $y$ is three stator slots, which leads to the angular coil width:
    \begin{equation}
        w = \frac{2\pi}{Q}m
        = \frac{2\pi}{15} \cdot 3
        = \frac{2\pi}{5}.
    \end{equation}

    Therefore, the chording factor is determined with:
    \begin{equation}
        s = \frac{w}{\tau_{\mathrm{p}}}
        = \frac{6\pi \cdot 2}{15 \pi}
        =\frac{12}{15}.
    \end{equation}

\end{solutionblock}

%%%%%%%%%%%%%%%%%%%%%%%%%%%%%%%%%%%%%%%%%%%%%%%%%%%%%%%%%%%%%
\subtask{Complete the winding scheme of the given machine in \autoref{tab:distributed_winding} for the phases b and c. Sketch the resulting winding scheme into \autoref{fig:MMF_distributed}.}

\begin{table}[ht]
    \caption{Winding scheme of the distributed winding from \autoref{fig:MMF_distributed}.}
    \centering
    \begin{tabular}{c|C{1cm}|C{1cm}|C{1cm}|C{1cm}|C{1cm}|C{1cm}}\toprule
        \multirow{2}{*}{Coil nr.} & \multicolumn{2}{c}{Phase a} & \multicolumn{2}{c}{Phase b} & \multicolumn{2}{c}{Phase c} \\
          & In  & Out   & In & Out & In & Out \\
        \midrule
        1  & 1  & 4  & & & & \\
        2  & 8  & 5  & & & & \\
        3  & 8  & 11 & & & & \\
        4  & 9  & 12 & & & & \\
        5  & 15 & 12 & & & & \\
        \bottomrule
    \end{tabular}
    \label{tab:distributed_winding}
\end{table}

\begin{solutionblock}
    In the task only phase a of the winding scheme is given. However, to complete the winding scheme, phase b must have an electrical phase shift of $\SI{120}{\degree}$ and phase c of $\SI{240}{\degree}$ respectively.
    This leads to
    \begin{equation}
        n_{\mathrm{b}} p \frac{\SI{360}{\degree}}{Q} = \SI{120}{\degree} \ \mathrm{mod} \ \SI{360}{\degree},
    \end{equation}
    where $n_{\mathrm{b}}$ is the displacement in slots. This expression is rewritten into
    \begin{equation}
        n_{\mathrm{b}} p \frac{\SI{360}{\degree}}{Q} = \SI{120}{\degree} + k \SI{360}{\degree},
        \label{eq:nb}
    \end{equation}
    and also for phase c:
    \begin{equation}
        n_{\mathrm{c}} p \frac{\SI{360}{\degree}}{Q} = \SI{240}{\degree} + k \SI{360}{\degree}.
        \label{eq:nc}
    \end{equation}
    
    Now, an integral solution must be found for \eqref{eq:nb} and \eqref{eq:nc} to determine the winding scheme.
    This is done with trial and error, the first assumption is $n_{\mathrm{b}} = 10, k = 1$, which results in
    \begin{align}
        \begin{split}
            10\cdot 2\cdot \frac{\SI{360}{\degree}}{15} &= \SI{120}{\degree}+\SI{360}{\degree}, \\
            \SI{480}{\degree} &= \SI{480}{\degree},
        \end{split}
    \end{align}
    therefore, the requirement is fulfilled. For phase c the assumption is given with $n_{\mathrm{c}} = 5, k = 0$ into \eqref{eq:nc}
    \begin{align}
        \begin{split}
            5 \cdot 2 \cdot \frac{\SI{360}{\degree}}{15} &= \SI{240}{\degree} + k \SI{360}{\degree}, \\
            \SI{240}{\degree} &= \SI{240}{\degree},
        \end{split}
    \end{align}
    which fulfills the requirement too.
    With the calculated phase shift and the knowledge of the coil width of three slots from the figure, the winding scheme is completed.
    In \autoref{tab:solution_distributedWinding} the inputs and outputs of each winding are given.
    \begin{solutiontable}[ht]
        \caption{Solution of the distributed winding scheme.}
        \centering
        \begin{tabular}{c|C{1cm}|C{1cm}|C{1cm}|C{1cm}|C{1cm}|C{1cm}}\toprule
            \multirow{2}{*}{Coil nr.} & \multicolumn{2}{c}{Phase a} & \multicolumn{2}{c}{Phase b} & \multicolumn{2}{c}{Phase c} \\
              & In  & Out   & In & Out & In & Out \\
            \midrule
            1  & 1  & 4  & 11 & 14 & 6  & 9 \\
            2  & 8  & 5  & 3  & 15 & 13 & 10 \\
            3  & 8  & 11 & 3  & 6  & 13 & 1 \\
            4  & 9  & 12 & 4  & 7  & 14 & 2 \\
            5  & 15 & 12 & 10 & 7  & 5  & 2 \\
            \bottomrule
        \end{tabular}
        \label{tab:solution_distributedWinding}
    \end{solutiontable}

    The sketch of the distributed winding scheme is visualized in \autoref{fig:solution_MMF_distributed}.
    \begin{solutionfigure}[ht]
        \centering
        \includegraphics{ex04/solution_MMF_distributed.pdf}
        \caption{Solution of the distributed winding.}
        \label{fig:solution_MMF_distributed}
    \end{solutionfigure}

\end{solutionblock}


%%%%%%%%%%%%%%%%%%%%%%%%%%%%%%%%%%%%%%%%%%%%%%%%%%%%%%%%%%%%%
\subtask{How many layers are there in this winding scheme?}

\begin{solutionblock}
    The winding diagram has 2 winding layers, due to the two conductors per slot.
\end{solutionblock}



%%%%%%%%%%%%%%%%%%%%%%%%%%%%%%%%%%%%%%%%%%%%%%%%%%%%%%%%%%%%%
%% Task 2 %%
%%%%%%%%%%%%%%%%%%%%%%%%%%%%%%%%%%%%%%%%%%%%%%%%%%%%%%%%%%%%%

\task{Concentrated windings}
Consider the shown 10-pole 3-phase motor with 12 stator slots in \autoref{fig:MMF_concentrated}. The winding scheme of phase a is shown in the figure.

\begin{figure}[htb]
    \centering
    \includegraphics{ex04/MMF_concentrated.pdf}
    %\captionsetup{labelfont={color=blue},textfont={color=blue}}
    \caption{Simplified sketch of a concentrated winding. Only phase a is shown.}
    \label{fig:MMF_concentrated}
\end{figure}

%%%%%%%%%%%%%%%%%%%%%%%%%%%%%%%%%%%%%%%%%%%%%%%%%%%%%%%%%%%%%
\subtask{Determine the pole pitch and und the number of notches. Is it an integral-slot or a fractional-slot winding?}

\begin{solutionblock}
    The number of slots is determined by counting the slots in the given figure of the machine.
    Hence, the number of slots is $Q$ = 12. The number of phases is given in the task, which is $m$=3.
    Therefore, the pole pitch calculates as:
    \begin{equation}
        \rho_{\mathrm{p}} = \frac{2\pi}{2p}
        = \frac{\pi}{5}.
    \end{equation}

    The number of notches is determined with:
    \begin{equation}
        q = \frac{Q}{2pm}
        = \frac{12}{2 \cdot 5 \cdot 3}
        = \frac{2}{5}.
    \end{equation}
    Since the number of notches per pole and phase is not an integer, it is a fractional slot winding.
    The coil width $y$ is one stator slot, as it is shown in the figure. This results in the angular coil width: 
    \begin{equation}
        w = \frac{2\pi}{Q}
        = \frac{2\pi}{12}
        = \frac{\pi}{6}.
    \end{equation}

    Hence, the chording factor is calculated by:
    \begin{equation}
        s = \frac{w}{\tau_{\mathrm{p}}}
        = \frac{\pi \cdot 5}{6\pi}
        = \frac{5}{6}.
    \end{equation}

\end{solutionblock}

%%%%%%%%%%%%%%%%%%%%%%%%%%%%%%%%%%%%%%%%%%%%%%%%%%%%%%%%%%%%%
\subtask{Complete \autoref{tab:concentrated_winding} with the winding scheme for phases b and c. Sketch the resulting winding scheme in \autoref{fig:MMF_concentrated}.}

\begin{table}[ht]
    \caption{Winding scheme of a concentrated winding from \autoref{fig:MMF_concentrated}.}
    \centering
    \begin{tabular}{c|C{1cm}|C{1cm}|C{1cm}|C{1cm}|C{1cm}|C{1cm}}\toprule
        \multirow{2}{*}{Coil nr.} & \multicolumn{2}{c}{Phase a} & \multicolumn{2}{c}{Phase b} & \multicolumn{2}{c}{Phase c} \\
          & In  & Out   & In & Out & In & Out \\
        \midrule
        1  & 2  & 1  & & & & \\
        2  & 12  & 1  & & & & \\
        3  & 7  & 6  & & & & \\
        4  & 7  & 8  & & & & \\
        \bottomrule
    \end{tabular}
    \label{tab:concentrated_winding}
\end{table}


\begin{solutionblock}
    In the task only phase a of the winding scheme is given. However, to complete the winding scheme, phase b must have an electrical phase shift of $\SI{120}{\degree}$ and phase c of $\SI{240}{\degree}$ respectively.
    This leads to
    \begin{equation}
        n_{\mathrm{b}} p \frac{\SI{360}{\degree}}{Q} = \SI{120}{\degree} \ \mathrm{mod} \ \SI{360}{\degree},
    \end{equation}
    where $n_{\mathrm{b}}$ is the displacement in slots. This expression is rewritten into
    \begin{equation}
        n_{\mathrm{b}} p \frac{\SI{360}{\degree}}{Q} = \SI{120}{\degree} + k \SI{360}{\degree},
        \label{eq:nb_concentrated}
    \end{equation}
    and also for phase c:
    \begin{equation}
        n_{\mathrm{c}} p \frac{\SI{360}{\degree}}{Q} = \SI{240}{\degree} + k \SI{360}{\degree}.
        \label{eq:nc_concentrated}
    \end{equation}
    
    Now, an integral solution must be found for \eqref{eq:nb_concentrated} and \eqref{eq:nc_concentrated} to determine the winding scheme.
    This is done with trial and error, the first assumption is $n_{\mathrm{b}} = 8, k = 3$, which results in
    \begin{align}
        \begin{split}
            8\cdot 5\cdot \frac{\SI{360}{\degree}}{12} &= \SI{120}{\degree}+3 \cdot \SI{360}{\degree}, \\
            \SI{1200}{\degree} &= \SI{1200}{\degree},
        \end{split}
    \end{align}
    therefore, the requirement is fulfilled. For phase c the assumption is given with $n_{\mathrm{c}} = 4, k = 1$ into \eqref{eq:nc_concentrated} leads to
    \begin{align}
        \begin{split}
            4 \cdot 5 \cdot \frac{\SI{360}{\degree}}{12} &= \SI{240}{\degree} + k \SI{360}{\degree}, \\
            \SI{600}{\degree} &= \SI{600}{\degree},
        \end{split}
    \end{align}
    which fulfills the requirement too.

    With the coil width of one slot, the resulting winding is determined. The solution of this winding scheme is shown in \autoref{tab:solution_concentratedWinding}.
    \begin{solutiontable}[ht]
        \caption{Solution of the winding scheme of a concentrated winding.}
        \centering
        \begin{tabular}{c|C{1cm}|C{1cm}|C{1cm}|C{1cm}|C{1cm}|C{1cm}}\toprule
            \multirow{2}{*}{Coil nr.} & \multicolumn{2}{c}{Phase a} & \multicolumn{2}{c}{Phase b} & \multicolumn{2}{c}{Phase c} \\
              & In  & Out   & In & Out & In & Out \\
            \midrule
            1  & 2  & 1  & 3 & 2 & 4 & 5 \\
            2  & 12 & 1  & 3 & 4 & 6 & 5 \\
            3  & 7  & 6  & 8 & 9 & 11 & 10 \\
            4  & 7  & 8  & 10 & 9 & 11 & 12\\
            \bottomrule
        \end{tabular}
        \label{tab:solution_concentratedWinding}
    \end{solutiontable}
    

    The completed winding scheme is visualized in Fig.\ref{fig:solution_MMF_concentrated}.
    \begin{solutionfigure}[ht]
        \centering
        \includegraphics{ex04/solution_MMF_concentrated.pdf}
        \captionsetup{labelfont={color=blue},textfont={color=blue}}
        \caption{Solution of the concentrated winding scheme.}
        \label{fig:solution_MMF_concentrated}
    \end{solutionfigure}
\end{solutionblock}



%%%%%%%%%%%%%%%%%%%%%%%%%%%%%%%%%%%%%%%%%%%%%%%%%%%%%%%%%%%%%
\subtask{How many layers are in this winding scheme?}

\begin{solutionblock}
    There are two layers per slot, therefore, it is a two layer winding scheme.
\end{solutionblock}